\documentclass[10pt,a4paper]{amsart}
\usepackage[margin=19mm]{geometry}
\usepackage{amsmath,amssymb,mathtools}
\usepackage[scr=rsfs,cal=euler]{mathalfa}
\usepackage{xfrac}
\usepackage[unicode,hidelinks,bookmarks=false]{hyperref}
\newcommand{\ie}{i.e.\ }
\newcommand{\cf}{cf.\ }
\newcommand{\resp}{respectively\ }
\newcommand{\Subsection}{\subsection}
\providecommand{\nobreakdash}{\nobreak\hbox{-}}
\setlength{\emergencystretch}{2em}
\multlinegap0pt
\usepackage{mathtools}
\newtheorem{lemma}{Lemma}
\newtheorem{theorem}{Theorem}
\usepackage{cleveref}
\crefname{lemma}{Lemma}{Lemmas}
\crefname{theorem}{Theorem}{Theorems}
\newcommand{\BB}{\mathcal B}
\newcommand{\DD}{\mathsf D}
\newcommand{\e}{\mathrm e}
\title{Retained-Set Descent for Diagonal Ramsey Numbers}
\author{Zhipeng Lu, Sichen Wang}
\date{Source: arXiv:2609.14525v1 (CC BY 4.0)}
\begin{document}
\maketitle

\noindent\emph{Source and licence.} Retained-Set Descent for Diagonal Ramsey Numbers. Version arXiv:2609.14525v1 (CC BY 4.0), distributed by arXiv under the Creative Commons Attribution 4.0 International licence, http://creativecommons.org/licenses/by/4.0/, as recorded in the arXiv licence metadata for this exact version and shown on the abstract page https://arxiv.org/abs/2609.14525v1. Full licence notice: \texttt{licenses/arxiv-2609.14525-CC-BY-4.0.txt}.\par\smallskip

\noindent\textit{Editorial scope.} Counted unit: the article's theorem thm:route (source0.tex lines 261--276). The article's complete original proof of it is delivered with it (source0.tex lines 278--348, one \texttt{proof} environment of 71 lines). No mathematics is written, repaired or re-derived: the statement and the proof are the article's own lines, in the article's own notation, and no retained line is substituted or re-flowed.  Retained uncounted as the source-local context the proof needs: lines 197--208.  Nothing was omitted from any retained span, so the package carries no omission record.  No retained line contains a citation, so the wrapper carries no bibliography and no bibliography entry.  No retained line contains a figure environment, an image-inclusion command or drawing code, so no image is delivered and the package declares no asset dependency.  Editorial formatting only, in addition: an AMS \texttt{amsart} preamble that restates the article's own declarations and macros used by the retained lines, namely the environments \texttt{lemma}, \texttt{theorem} on separate counters, together with the standard packages those lines need.  The retained environments print in this document as Lemma 1, Theorem 1 in order of appearance, whereas the article prints the same items with the numbering of its own full document.  The complete native source of the article as distributed in the arXiv e-print for this exact version is bundled unchanged as \texttt{sources/210327/source0.tex}.
\begin{lemma}[Weighted Candidate Bound]\label{prop:candidate}
Fix $w>0$ and $\mu,p,x,y\in(0,1)$ with $(x,y)\in\operatorname{int}\BB_f$ and
\[
 x<(1-\mu)^w p^{1/(1-\mu)}.
\]
For all sufficiently large $\ell$, uniformly in positive integers $k,t$,
\[
 d_R(X,Y)\ge p,\qquad
 |X|^w|Y|\ge x^{-k}y^{-\ell}\mu^{-wt}
\]
imply that $(X,Y)$ is $(k,\ell,t)$-good.
\end{lemma}

\begin{theorem}[Retained-Set Rule]\label{thm:route}
Fix a route as above, an outer density $p$, and a row density $0<\pi<p<1$.
Choose a strict control at density $\pi$ as the route's \emph{terminal parameters}.
If
\begin{equation}\label{eq:route-test}
 z>g(\lambda),\qquad
 z+w g(\theta)>A+\lambda B+w\theta\beta,
\end{equation}
then $\DD(p,z,\{\lambda\})$ holds with
\begin{equation}\label{eq:C}
 C=\frac{3(1-\pi)}{p-\pi}.
\end{equation}
The conclusion is uniform over truncations and vertical shifts of one fixed route when the new endpoint, output exponent, and shift depend continuously on a parameter in a compact set; the stopping ratio, slopes, densities, calls, and terminal controls stay fixed.
At endpoint $r$, retain the closed cells $[u_{j-1},\min\{u_j,r\}]$ with $u_{j-1}<r$ and require positive uniform margins in $g>L_j$ and in~\Cref{eq:route-test}.
Such a family may include $r=\theta$, with no internal calls.
\end{theorem}

\begin{proof}
If $\lambda=\theta$, the size tests give
\[
 (1+w)z>z+wg(\theta)>A+\theta B+w\theta\beta.
\]
The source leaf at density $\pi$ proves the conclusion at density $p>\pi$, with constant $1\le C$.

Suppose $\theta<\lambda$, and take a graph of red density at least $p$ and order $N\ge C\e^{kz}$.
Averaging gives a balanced partition $X_0,Y$ with cross-density at least $p$ and both parts of size at least $N/3$.
Keep in $X_0$ only the vertices with at least $\pi|Y|$ red neighbors in $Y$.
The surviving set $X$ satisfies
\[
 (p-\pi)|X_0||Y|
 \le e_R(X,Y)-\pi|X||Y|
 \le(1-\pi)|X||Y|.
\]
The choice of $C$ gives $|X|,|Y|\ge\e^{kz}$.
Every nonempty subset of $X$ has red cross-density at least $\pi$ with this fixed $Y$.

Put $n_0=\lfloor\theta k\rfloor$ and $\ell=\lfloor\lambda k\rfloor$, taking $k$ large enough that $n_0\ge1$.
To track the loss, define for integers $0\le t\le\ell$
\[
 \Gamma_k(t)=kg(\theta)+\sum_{j=n_0+1}^{t}d(j/k)\quad(t>n_0),
 \qquad \Gamma_k(t)=kg(\theta)\quad(0\le t\le n_0).
\]
For $V=d_1+\sum_{j<m}|d_{j+1}-d_j|$,
\begin{equation}\label{eq:discrete}
 \Gamma_k(t)-\Gamma_k(t-1)=d(t/k)\quad(t>n_0),\qquad
 |\Gamma_k(t)-kg(t/k)|\le V.
\end{equation}
For the second bound, express $d$ as its initial step at $\theta$ and its subsequent jumps.
For each unit step, the mesh count differs from $k$ times its integral by at most one, even when a jump lies on the mesh.
Summing the absolute jump sizes gives $V$.

Let $C_j$ be the call constants and put
\[
 \delta=\min_j\min_{J_j}(g-L_j)>0,\qquad
 \eta=\min_j(1-q_j-\e^{-d_j})>0.
\]
Choose $k$ so that $k\delta>V+\max_j\log C_j$ and $\eta\e^{kg(\theta)}>1$, and so that all call and candidate thresholds hold.
We prove by induction on $t=1,\ldots,\ell$ that $(Z,Y)$ is $(k,\ell,t)$-good whenever $Z\subseteq X$ and $|Z|\ge\e^{\Gamma_k(t)}$.

For $t\le n_0$, the terminal test gives
\[
 \log(|Z|^w|Y|)\ge k[wg(\theta)+z]
 >k(A+\lambda B+w\theta\beta)
 \ge kA+\ell B+wt\beta.
\]
Since $d_R(Z,Y)\ge\pi$, \Cref{prop:candidate} applies.

For $t>n_0$, put $s=t/k$ and use the cell with assigned slope $d_j$.
If the red density in $Z$ is at least $q_j$, then
$\log|Z|\ge kg(s)-V>kL_j(s)+\log C_j$.
The call gives a red $K_k$ or a blue $K_t$.
Otherwise some vertex has a blue neighborhood $Z'\subseteq Z$ with
\[
 |Z'|>(\e^{-d_j}+\eta)(|Z|-1)>\e^{-d_j}|Z|
 \ge\e^{\Gamma_k(t-1)}.
\]
The middle inequality follows from $\eta|Z|>1$ and $\e^{-d_j}+\eta<1$.
Induction applies to $(Z',Y)$.
A blue $K_{t-1}$ extends through the chosen vertex; either of the other two outcomes already proves the claim.

Finally, \Cref{eq:discrete} and the endpoint test give
$\Gamma_k(\ell)\le kg(\lambda)+V<kz$ for large $k$.
The claim therefore applies to $Z=X$ at $t=\ell$.

For uniformity, take the minimum stopping height and the minima of the strict gaps over the compact parameter set, and the maxima of the finitely many call thresholds.
The same $V$ bounds every truncated profile, since truncation only removes slope jumps.
These choices give a common threshold for $k$; at $\lambda=\theta$ use the fixed source control as above.
\end{proof}

\end{document}
